\begin{figure}[tbp]
\centering
\begin{adjustbox}{max width=\textwidth,center}
\begin{tikzpicture}[>=Stealth,
 vertex/.style={draw,rounded corners,fill=blue!3,
 minimum width=2.5cm,minimum height=.9cm,align=center},
 edge/.style={->,thick,blue!60!black}]
\node[vertex] (one) at (0,2.2) {$F_1$\\centered curvature $\tau$};
\node[vertex] (three) at (-2.8,0) {$F_3$\\centered curvature $0$};
\node[vertex] (two) at (2.8,0) {$F_2$\\centered curvature $-2\tau/5$};
\node[vertex] (four) at (0,-2.2) {$F_4$\\centered curvature $-3\tau/5$};
\draw[edge] (two.north) |- (one.east)
 node[pos=.28,right] {$x_{12}$};
\draw[edge] (three.north) |- (one.west)
 node[pos=.28,left] {$x_{13}$};
\draw[edge] (two.west) -- (three.east)
 node[midway,above] {$x_{32}$};
\draw[edge] (four.east) -| (two.south)
 node[pos=.70,right] {$x_{24}$};
\draw[edge] (four.west) -| (three.south)
 node[pos=.70,left] {$x_{34}$};
\end{tikzpicture}
\end{adjustbox}
\caption{The five allowed harmonic blocks on the unstable side.
Arrows run from source to target. For $\tau>0$ their line
quotients have the decreasing slope order $1,3,2,4$.
The full reconstruction also contains induced nonlinear entries;
the diagram does not replace the integrability equation $\kappa=0$.}
\label{fig:unstable-support}
\end{figure}
